% ============ 10 \v{C}ech 上同调 (原书 pp.70-75) ============
\chapter{\v{C}ech 上同调}\label{sec:10}
\index{\v{C}ech 上同调}
\index{Mayer-Vietoris 序列}
\index{谱序列}
\index{逆向极限}

用导出函子的定义直接计算上同调群是不实际的.在 $X$ 的温和假设下，导出函子群与 \v{C}ech 群一致，后者有时更易处理.

\subsection*{\v{C}ech 群的定义}

设 $\mathcal{U} = (U_i \to X)_{i \in I}$ 是 $X$ 的平展覆盖，$P$ 是 $X_{\mathrm{et}}$ 上的 Abel 群预层.定义
\[
C^{r}(\mathcal{U}, P) = \prod_{(i_0, \ldots, i_r) \in I^{r+1}} P(U_{i_0 \cdots i_r}), \qquad \text{其中 } U_{i_0 \cdots i_r} = U_{i_0} \times_{X} \cdots \times_{X} U_{i_r}.
\]
对 $s = (s_{i_0 \cdots i_r}) \in C^{r}(\mathcal{U}, P)$，按规则定义 $d^{r}s \in C^{r+1}(\mathcal{U}, P)$：
\[
(d^{r}s)_{i_0 \cdots i_{r+1}} = \sum_{j=0}^{r+1} (-1)^{j} \operatorname{res}_j(s_{i_0 \cdots i_{j-1} i_{j+1} \cdots i_{r+1}}),
\]
其中 $\operatorname{res}_j$ 是对应于投影映射
\[
U_{i_0 \cdots i_{r+1}} \to U_{i_0 \cdots i_{j-1} i_{j+1} \cdots i_{r+1}}
\]
的限制映射.

与经典情形一样，通过直接计算验证
\[
C^{\bullet}(\mathcal{U}, P) \overset{\mathrm{def}}{=} C^{0}(\mathcal{U}, P) \to \cdots \to C^{r}(\mathcal{U}, P) \xrightarrow{d^{r}} C^{r+1}(\mathcal{U}, P) \to \cdots
\]
是复形.定义
\[
\check{H}^{r}(\mathcal{U}, P) = H^{r}(C^{\bullet}(\mathcal{U}, P)).
\]
它称为 $P$ 关于覆盖 $\mathcal{U}$ 的第 $r$ 个 \v{C}ech 上同调群.

注意
\[
\check{H}^{0}(\mathcal{U}, P) = \operatorname{Ker}\left( \prod P(U_i) \rightrightarrows \prod P(U_{ij}) \right).
\]
因此，对层 $F$，
\[
\check{H}^{0}(\mathcal{U}, F) = \Gamma(X, F).
\]

\begin{example}\label{ex:10.1}
设 $\mathcal{U}$ 是由带 Galois 群 $G$ 的单个 Galois 覆盖 $Y \to X$ 组成的 $X$ 的覆盖.若 $P$ 是 $X$ 上把不交并送到乘积的预层，则
\[
\check{H}^{r}(\mathcal{U}, P) = H^{r}(G, P(Y)) \quad (\text{群上同调});
\]
见习题~\ref{exc:6.5}.
\end{example}

$X$ 的第二个覆盖 $\mathcal{V} = (V_j \to X)_{j \in J}$ 称为 $\mathcal{U}$ 的加细，如果存在映射 $\psi \colon J \to I$，使得对所有 $j \in J$，$V_j \to X$ 穿过 $U_{\psi(j)} \to X$.$\psi$ 与 $X$-态射 $\varphi_j \colon V_j \to U_{\psi(j)}$(对每个 $j$)的选取确定复形映射
\[
\psi^{\bullet} \colon C^{\bullet}(\mathcal{U}, P) \to C^{\bullet}(\mathcal{V}, P), \qquad (\psi^{r}s)_{j_0 \cdots j_r} = s_{\psi(j_0) \cdots \psi(j_r)}|_{V_{j_0 \cdots j_r}}.
\]
与经典情形一样，可以验证上同调群上的映射
\[
\rho(\mathcal{V}, \mathcal{U}) \colon \check{H}^{r}(\mathcal{U}, P) \to \check{H}^{r}(\mathcal{V}, P)
\]
不依赖于所有选择.我们可以对所有覆盖取极限，从而得到 \v{C}ech 上同调群
\[
\check{H}^{r}(X, P) \overset{\mathrm{def}}{=} \varinjlim_{\mathcal{U}} \check{H}^{r}(\mathcal{U}, P).
\]
它们有如下性质：
\begin{enumerate}
  \item 对 $X$ 上的任意层 $F$，$\check{H}^{0}(X, F) = \Gamma(X, F)$；
  \item 对所有内射层 $\mathcal{I}$, $\check{H}^{r}(X, \mathcal{I}) = 0$($r > 0$).
\end{enumerate}
论断 (a) 显然.论断 (b) 通过证明如下事实来证：对每个覆盖 $(U_i \to X)$，存在 Abel 群预层的正合序列 $Z^{\bullet}$，使得对所有层 $F$，
\[
C^{\bullet}(\mathcal{U}, F) = \operatorname{Hom}(Z^{\bullet}, F).
\]
若 $\mathcal{I}$ 作为层内射，则它作为预层也内射(因为 $a$ 是包含函子的正合左伴随；比较 7.3c)，从而 $\operatorname{Hom}(Z^{\bullet}, \mathcal{I})$ 正合.细节见 EC III 2.4 的证明.

由此推出，同构 $\check{H}^{0}(X, F) \cong H^{0}(X, F)$ 对所有 $r$ 与 $F$ 延拓为同构，当且仅当层的每个短正合序列
\[
0 \to F' \to F \to F'' \to 0
\]
都给出 \v{C}ech 上同调群的长正合序列
\[
\cdots \to \check{H}^{r}(X, F') \to \check{H}^{r}(X, F) \to \check{H}^{r}(X, F'') \to \cdots.
\]

\begin{theorem}\label{thm:10.2}
假定 $X$ 的每个有限子集都包含在开仿射集中，且 $X$ 拟紧(例如 $X$ 可以是拟射影簇).则对层的任何短正合序列
\[
0 \to F' \to F \to F'' \to 0,
\]
复形
\[
0 \to C^{\bullet}(\mathcal{U}, F') \to C^{\bullet}(\mathcal{U}, F) \to C^{\bullet}(\mathcal{U}, F'') \to 0
\]
在 $X$ 的平展覆盖上的直限正合，从而给出 \v{C}ech 上同调群的长正合序列.因此对所有 $r$ 与所有层 $F$，
\[
\check{H}^{r}(X, F) \cong H^{r}(X, F).
\]
\end{theorem}

证明 \v{C}ech 复形的直限的正合性的困难如下：由于 $F \to F'' \to 0$ 正合，我们知道 $F(U_{i_0 \cdots i_n}) \to F''(U_{i_0 \cdots i_n})$ 局部满，即对每个 $s \in F''(U_{i_0 \cdots i_n})$，存在覆盖 $(V_j \to U_{i_0 \cdots i_n})_{j \in J}$，使得对每个 $j$，$s|_{V_j}$ 提升到 $F(V_j)$；问题是一般地我们不知道诸 $V_j$ 能否取为形式 $V_{i_0 \cdots i_n} = V_{i_0} \times_{X} \cdots \times_{X} V_{i_n}$，其中 $V_{i_m}$ 映到 $U_{i_m}$.

定理证明的关键是 M. Artin 的如下结果(Advances in Math. 7 (1971), 282--296.)(另见 Hochster 和 Huneke, Ann. of Math., 135 (1992)，定理 9.2):

\begin{quote}
设 $A$ 是环，$p_1, \ldots, p_r$ 是 $A$ 中的素理想，$A_1, \ldots, A_r$ 是局部环 $A_{p_1}, \ldots, A_{p_r}$ 的严格 Hensel 化.则 $A_0 \overset{\mathrm{def}}{=} A_1 \otimes_{A} \cdots \otimes_{A} A_r$ 有性质：任何忠实平坦平展映射 $A_0 \to B$ 都有截面 $B \to A_0$.
\end{quote}

几乎按定义，严格局部环有这个性质.这个论断的趣味在于它对从单个环得到的严格局部环的张量积成立.

定理~\ref{thm:10.2} 证明的其余部分见 Artin 的论文或 EC III 2.17.

\begin{remark}\label{rem:10.3}
(a) 对 Zariski 拓扑，只有较弱的结论为真：对簇或分离概形，凝聚 $\mathcal{O}_X$-模的 \v{C}ech 上同调与导出函子上同调一致.这由如下事实推出：对任何开仿射 $U$ 与凝聚 $\mathcal{O}_X$-模的正合序列
\[
0 \to \mathcal{M}' \to \mathcal{M} \to \mathcal{M}'' \to 0,
\]
\[
0 \to \Gamma(U, \mathcal{M}') \to \Gamma(U, \mathcal{M}) \to \Gamma(U, \mathcal{M}'') \to 0
\]
正合.因此，对应于 $X$ 的开仿射覆盖 $X = \bigcup U_i$ 的 \v{C}ech 复形正合(因为分离空间中开仿射的有限交是仿射的---AG 4.27).

当然，常值层的 \v{C}ech 上同调群与导出函子上同调群一致：当 $X$ 不可约时它们都为零.

(b) 关于拓扑空间的 \v{C}ech 上同调的大部分形式体系都适用于平展景，但不是全部.例如，\v{C}ech 上同调不能用交错上链计算，如例~\ref{ex:10.1} 所示.通常的证明失效，因为从一个平展 $X$-簇到另一个可能有多个映射.
\end{remark}

\subsection*{第一上同调群的比较}

在任何景上，第一 \v{C}ech 上同调群等于第一导出函子群.我们概述证明.

\begin{proposition}\label{prop:10.4}
对 $X_{\mathrm{et}}$ 上的层 $F$，设 $\mathcal{H}^{r}(F)$ 是预层 $U \mapsto H^{r}(U, F|_{U})$.对所有 $r > 0$，与 $\mathcal{H}^{r}(F)$ 相伴的层是 0.
\end{proposition}

\begin{proof}
考虑函子
\[
\mathrm{Sh}(X_{\mathrm{et}}) \xrightarrow{i} \mathrm{PreSh}(X_{\mathrm{et}}) \xrightarrow{a} \mathrm{Sh}(X_{\mathrm{et}}).
\]
回忆 \S 7 中 $i$ 左正合且 $a$ 正合.设 $F \to \mathcal{I}^{\bullet}$ 是 $F$ 的内射分解.则
\[
\mathcal{H}^{r}(F) = H^{r}(i\mathcal{I}^{\bullet}).
\]
因为 $a$ 正合且 $a \circ i = \mathrm{id}$，
\[
a(\mathcal{H}^{r}(F)) \overset{\mathrm{def}}{=} a(H^{r}(i\mathcal{I}^{\bullet})) = H^{r}(ai\mathcal{I}^{\bullet}) = H^{r}(\mathcal{I}^{\bullet}) = 0 \quad (r > 0).
\]
\end{proof}

\begin{corollary}\label{cor:10.5}
设 $s \in H^{r}(X, F)$($r > 0$ 适当).则存在覆盖 $(U_i \to X)$，使得 $s$ 在每个群 $H^{r}(U_i, F)$ 中的像为零.
\end{corollary}

\begin{proof}
回忆对任何预层 $P$，$P$ 中映到 $aP$ 的零的截面恰好是局部为零的截面，即限制到 $X$ 的某个覆盖中的 $U_i$ 上为零的截面.
\end{proof}

现在我们定义映射 $H^{1}(X, F) \to \check{H}^{1}(X, F)$.选取 $F$ 的内射嵌入 $F \hookrightarrow \mathcal{I}$，设 $\mathcal{G}$ 是余核.于是，有正合序列
\[
0 \to F \to \mathcal{I} \to \mathcal{G} \to 0,
\]
它产生正合上同调序列
\[
0 \to F(X) \to \mathcal{I}(X) \to \mathcal{G}(X) \to H^{1}(X, F) \to 0.
\]
设 $s \in H^{1}(X, F)$，$t \in \mathcal{G}(X)$ 映到 $s$.由上面的推论，存在覆盖 $(U_i \to X)$，使得 $s$ 在每个 $U_i$ 上限制为零，从而 $t|_{U_i}$ 提升为 $\mathcal{I}(U_i)$ 的元素 $\tilde{t}_i$.设 $s_{ij} = \tilde{t}_j|_{U_{ij}} - \tilde{t}_i|_{U_{ij}}$(看作 $F(U_{ij})$ 的元素).容易检验 $s_{ij}$ 是一个 1-上链.

\begin{proposition}\label{prop:10.6}
映射 $s \mapsto (s_{ij})$ 定义同构 $H^{1}(X, F) \to \check{H}^{1}(X, F)$.
\end{proposition}

\begin{proof}
我们留作习题让读者找直接证明.或者，我们下面给出一个涉及谱序列的证明.
\end{proof}

\subsection*{联系 \v{C}ech 上同调与导出函子上同调的谱序列}

由于我们将多次使用它，我陈述 Grothendieck 关于谱序列存在性的定理.

\begin{theorem}\label{thm:10.7}
设 $A$，$B$，$C$ 是 Abel 范畴，假定 $A$ 与 $B$ 有足够的内射对象.设 $F \colon A \to B$ 与 $G \colon B \to C$ 是左正合函子，并假定若 $\mathcal{I}$ 内射则对 $r > 0$，$(R^{r}G)(F\mathcal{I}) = 0$(例如当 $F$ 把内射对象送到内射对象时成立).则存在谱序列
\[
E^{rs}_2 = (R^{r}G)(R^{s}F)(A) \Rightarrow R^{r+s}(FG)(A).
\]
\end{theorem}

EC pp.307--309 有三页的谱序列解释，Shatz, \emph{Profinite Groups, Arithmetic, and Geometry}, Princeton, 1972, II.4 有十四页的解释，Weibel 第 V 章有四十五页的解释.

现在我们概述联系 \v{C}ech 上同调与导出函子上同调的谱序列的推导.

可以证明 $X_{\mathrm{et}}$ 上的预层范畴有足够的内射对象.与层情形的论证相同，对内射预层 $\mathcal{I}$，对 $r > 0$，$\check{H}^{r}(X_{\mathrm{et}}, \mathcal{I}) = 0$.由于预层的短正合序列显然给出 \v{C}ech 上同调群的长正合序列，我们看到 $\check{H}^{r}(X_{\mathrm{et}}, \cdot)$ 是
\[
P \mapsto \check{H}^{0}(X_{\mathrm{et}}, P) \colon \mathrm{PreSh}(X_{\mathrm{et}}) \to \mathbf{Ab}
\]
的第 $r$ 个右导出函子.

考虑函子序列
\[
\mathrm{Sh}(X_{\mathrm{et}}) \xrightarrow{i} \mathrm{PreSh}(X_{\mathrm{et}}) \xrightarrow{\check{H}^{0}(X_{\mathrm{et}}, \cdot)} \mathbf{Ab}.
\]
$i$ 的第 $r$ 个右导出函子是 $\mathcal{H}^{r}(\cdot)$.我们已经注意到 $i$ 保持内射对象，因此 Grothendieck 的一个定理(Weibel 1994, 5.8.3)为我们提供谱序列：
\[
\check{H}^{r}(X_{\mathrm{et}}, \mathcal{H}^{s}(F)) \Rightarrow H^{r+s}(X_{\mathrm{et}}, F).
\]
由推论~\ref{cor:10.5}，对 $s > 0$，$\check{H}^{0}(X_{\mathrm{et}}, \mathcal{H}^{s}(F)) = 0$.于是，对层 $F$，
\[
\check{H}^{r}(X, F) \cong H^{r}(X, F) \quad (r = 0, 1),
\]
且有正合序列
\[
0 \to \check{H}^{2}(X, F) \to H^{2}(X, F) \to \check{H}^{1}(X, \mathcal{H}^{1}(F)) \to \check{H}^{3}(X, F) \to H^{3}(X, F).
\]
类似地，对 $X$ 的任何平展覆盖 $\mathcal{U} = (U_i \to X)$，存在谱序列
\[
\check{H}^{r}(\mathcal{U}, \mathcal{H}^{s}(F)) \Rightarrow H^{r+s}(X_{\mathrm{et}}, F).
\]

\textbf{Mayer-Vietoris 序列} 当 $\mathcal{U} = (U_i \to X)$ 是 $X$ 的(Zariski 意义下的)开覆盖时，\v{C}ech 上同调群可以用交错上链计算.例如，若 $X = U_0 \cup U_1$，则预层 $P$ 的 \v{C}ech 上同调群是复形
\[
\Gamma(U_0, P) \oplus \Gamma(U_1, P) \to \Gamma(U_0 \cap U_1, P)
\]
的上同调群；特别地，对 $r \geq 2$，$\check{H}^{r}(\mathcal{U}, P) = 0$.

\begin{theorem}\label{thm:10.8}
设 $X = U_0 \cup U_1$(两个开子集的并).对 $X_{\mathrm{et}}$ 上的任意层 $F$，有无限正合序列
\[
\cdots \to H^{s}(X, F) \to H^{s}(U_0, F) \oplus H^{s}(U_1, F) \to H^{s}(U_0 \cap U_1, F) \to H^{s+1}(X, F) \to \cdots.
\]
\end{theorem}

\begin{proof}
在上述讨论中取 $P = \mathcal{H}^{s}(F)$($s > 0$)，得到正合序列
\[
0 \to \check{H}^{0}(\mathcal{U}, \mathcal{H}^{s}(F)) \to H^{s}(U_0, F) \oplus H^{s}(U_1, F) \to H^{s}(U_0 \cap U_1, F) \to \check{H}^{1}(\mathcal{U}, \mathcal{H}^{s}(F)) \to 0.
\]
在谱序列
\[
\check{H}^{r}(\mathcal{U}, \mathcal{H}^{s}(F)) \Rightarrow H^{r+s}(X, F)
\]
中，除非 $r = 0, 1$，$\check{H}^{r}(\mathcal{U}, \mathcal{H}^{s}(F)) = 0$.因此谱序列给出正合序列
\[
0 \to \check{H}^{1}(\mathcal{U}, \mathcal{H}^{s}(F)) \to H^{s+1}(X, F) \to \check{H}^{0}(\mathcal{U}, \mathcal{H}^{s+1}(F)) \to 0 \quad (\text{所有 } s \geq 0).
\]
把这些序列拼接起来，得到定理叙述中的序列.
\end{proof}

\subsection*{上同调对概形逆向极限的行为}

设 $I$ 是有向集，$(X_i)_{i \in I}$ 是以 $I$ 为指标集的概形的射影系统.若转移映射
\[
X_i \xleftarrow{\varphi^{j}_i} X_j, \qquad i \leq j,
\]
仿射，即对任何开仿射 $U \subset X_i$，$(\varphi^{j}_i)^{-1}(U)$ 仿射，则逆向极限概形 $X_\infty = \varprojlim X_i$ 存在.例如，若所有 $X_i$ 都仿射，设 $X_i = \operatorname{Spec}(A_i)$，则我们有环的正向系统 $(A_i)_{i \in I}$，且 $X_\infty = \operatorname{Spec} A_\infty$，其中 $A_\infty = \varinjlim A_i$.

\begin{theorem}\label{thm:10.9}
设 $I$ 是有向集，$(X_i)_{i \in I}$ 是 $X$-概形的逆向系统.假定所有 $X_i$ 都拟紧，且映射 $X_i \to X_j$ 都仿射.设 $X_\infty = \varprojlim X_i$，对 $X$ 上的任意层 $F$，设 $F_i$ 是它在 $X_i$ 上的逆像($i \in I \cup \{\infty\}$).则
\[
\varinjlim_{i} H^{r}(X_{i,\mathrm{et}}, F_i) \cong H^{r}(X_\infty, F_\infty).
\]
\end{theorem}

\begin{proof}
一般情形的证明相当困难.它依赖于如下事实：$X_\infty$ 上有限型的平展概形范畴是诸 $X_i$ 上这样的范畴的正向极限.见 SGA4, VII.5.8，或 Artin 1962, III.3.

然而，对 \v{C}ech 上同调证明定理要容易得多(见 EC III.3.17).当 10.2 适用时，可以恢复导出函子群的定理.
\end{proof}
